% paper/v2/preamble.tex
% LaTeX preamble loaded by Pandoc via --include-in-header.
% Target: academic article typography (amsart class) with full theorem-
% environment support and unicode glyph fallback for STIX Two Text.

% --- Mathematical typography ---------------------------------------------
\usepackage{amsmath}
\usepackage{amssymb}
\usepackage{amsthm}
\usepackage{mathtools}

% --- Diagrams (TikZ for §8 dependency graph) -----------------------------
\usepackage{tikz}
\usetikzlibrary{positioning, arrows.meta, shapes.geometric}

% --- Tables (booktabs for proper rules) ----------------------------------
\usepackage{booktabs}
\usepackage{array}

% --- Running heads ------------------------------------------------------
% Academic-paper convention: bare folio at page bottom, no top header.
% This avoids the amsart-default behaviour of stuffing title + full
% author block into the running header and the resulting overprinting.
\pagestyle{plain}

% --- Theorem-like environments (numbered within sections) ----------------
\theoremstyle{plain}
\newtheorem{theorem}{Theorem}[section]
\newtheorem{lemma}[theorem]{Lemma}
\newtheorem{corollary}[theorem]{Corollary}
\newtheorem{proposition}[theorem]{Proposition}

\theoremstyle{definition}
\newtheorem{definition}[theorem]{Definition}
\newtheorem{conjecture}[theorem]{Conjecture}
\newtheorem{reformulation}[theorem]{Reformulation}

\theoremstyle{remark}
\newtheorem*{remark}{Remark}
\newtheorem*{note}{Note}

% --- Unicode glyph fallback for STIX Two Text gaps -----------------------
% Some Italic / Bold-Italic shapes of STIX Two Text lack arrows and a few
% relations; route them through math mode (which uses STIX Two Math).
\usepackage{newunicodechar}

% Arrows
\newunicodechar{→}{\ensuremath{\rightarrow}}
\newunicodechar{←}{\ensuremath{\leftarrow}}
\newunicodechar{↔}{\ensuremath{\leftrightarrow}}
\newunicodechar{⇒}{\ensuremath{\Rightarrow}}
\newunicodechar{⇐}{\ensuremath{\Leftarrow}}
\newunicodechar{⇔}{\ensuremath{\Leftrightarrow}}

% Order / equality relations
\newunicodechar{≤}{\ensuremath{\leq}}
\newunicodechar{≥}{\ensuremath{\geq}}
\newunicodechar{≠}{\ensuremath{\neq}}
\newunicodechar{≈}{\ensuremath{\approx}}
\newunicodechar{≡}{\ensuremath{\equiv}}
\newunicodechar{∼}{\ensuremath{\sim}}
